\section{Deployment e Evolução}

A estratégia de deployment deve preservar a independência dos três repositórios e, ao mesmo tempo, permitir a execução integrada do sistema.

\subsection{Desenvolvimento local}

Durante o desenvolvimento, as imagens podem ser construídas localmente a partir dos três repositórios. Um Compose de integração deve conectar frontend, Java, pipeline e PostgreSQL em uma rede comum.

\subsection{Integração contínua}

Em uma evolução futura, cada repositório poderá possuir seu próprio pipeline de CI/CD:

\begin{lstlisting}
commit
  |
  v
CI -> testes -> build -> imagem Docker -> registry
\end{lstlisting}

O repositório de deployment poderá selecionar versões específicas das imagens para cada ambiente.

\subsection{Versionamento das imagens}

É recomendável utilizar tags imutáveis ou identificáveis, como versões semânticas ou SHA do commit. Um ambiente integrado pode então combinar, por exemplo, versões diferentes do frontend, Java e pipeline de maneira explícita.

\subsection{Evolução da equipe}

A separação por repositórios permite que equipes futuras assumam componentes diferentes. Para evitar que a separação gere acoplamento descontrolado, devem ser mantidos:

\begin{itemize}
    \item contratos documentados;
    \item migrations versionadas;
    \item testes de integração;
    \item convenções de versionamento;
    \item documentação centralizada neste repositório.
\end{itemize}

\subsection{Evoluções futuras}

A arquitetura pode evoluir gradualmente para mecanismos como registry privado, CI/CD, observabilidade centralizada, jobs agendados e ambientes de staging. Essas ferramentas devem ser adotadas quando houver necessidade operacional real, evitando complexidade prematura.

\begin{explicacao}[Princípio de evolução incremental]
A arquitetura atual deve resolver o problema presente sem impedir crescimento futuro. Novas ferramentas e serviços devem ser introduzidos quando reduzirem um problema concreto de operação, manutenção ou escala.
\end{explicacao}
